\documentclass[10pt,a4paper]{amsart}
\usepackage[margin=19mm]{geometry}
\usepackage{amsmath,amssymb,mathtools}
\usepackage[scr=rsfs,cal=euler]{mathalfa}
\usepackage{xfrac}
\usepackage[unicode,hidelinks,bookmarks=false]{hyperref}
\newcommand{\ie}{i.e.\ }
\newcommand{\cf}{cf.\ }
\newcommand{\resp}{respectively\ }
\newcommand{\Subsection}{\subsection}
\providecommand{\nobreakdash}{\nobreak\hbox{-}}
\setlength{\emergencystretch}{2em}
\multlinegap0pt
\usepackage{dsfont}
\usepackage{bbm}
\usepackage{enumitem}
\newcommand{\R}{{\mathbb{R}}}
\newcommand{\f}{\frac}
\newcommand{\diff}{\mathop{}\!\mathrm{d}}
\newcommand{\ve}{\varepsilon}
\newcommand{\ind}{\mathbbm{1}}
\newcommand{\impl}{\Rightarrow}
\theoremstyle{plain}
\newtheorem{theorem}{Theorem}[section]
\newtheorem{lemma}[theorem]{Lemma}
\newtheorem{proposition}[theorem]{Proposition}
\newtheorem{corollary}[theorem]{Corollary}
\newtheorem{conjecture}[theorem]{Conjecture}
\newtheorem{problem}[theorem]{Problem}
\theoremstyle{remark}
\newtheorem{remark}[theorem]{\bf Remark}
\newtheorem{definition}[theorem]{\bf Definition}
\newlist{enumeratesteps}{enumerate}{2}
\setlist[enumeratesteps,1]{
  label=\textit{Step \arabic*\protect\thisstep},
  ref=\arabic*,
  align=left, leftmargin=0pt, labelindent=!,
  listparindent=\parindent, labelwidth=0pt, itemindent=!
}
\setlist[enumeratesteps,2]{
    label= \textit{Step \theenumeratestepsi.\alph*\protect\thisstep},
    ref=\theenumeratestepsi.\alph*,
    align=left, leftmargin=0pt, labelindent=\parindent,
  listparindent=\parindent, labelwidth=0pt, itemindent=!
}
\newcommand{\step}[1][]{%
  \if\relax\detokenize{#1}\relax
    \def\thisstep{.}%
  \else
    \def\thisstep{: \protect #1.}%
  \fi
  \item}
\title{A Li--Yau and Aronson--B\'enilan approach for the Keller--Segel system with critical exponent}
\author{Charles Elbar, Alejandro Fernández-Jiménez and Filippo Santambrogio}
\date{Source: arXiv:2512.17772v2 (CC BY 4.0); the authors' own native \LaTeX{} source, set here in the editorial AMS \texttt{amsart} wrapper (the source's own class is \texttt{article}, which is not redistributed).}
\begin{document}
\maketitle
\noindent\textit{Editorial scope.} The counted claim of this unit is the article's Theorem 6.9 (source label \texttt{thm:new\_proof\_hls}), the authors' new proof of the logarithmic Hardy--Littlewood--Sobolev inequality for subcritical and critical mass, restated in the introduction as Theorem~1.11, delivered together with the authors' complete original proof. The statement and the proof are the retained mathematical passages, and every one of their characters is the authors' own. The three commented-out source lines that the authors left between the statement and the proof are not retained, because a body consisting only of comments carries no mathematical content and the delivery gate rejects an empty context body; they print nothing in any case. Printed numbering: the authors' own theorem-environment declarations are carried unchanged, including the numbering of theorem-like environments by section, and the section and theorem counters are advanced so that the delivered PDF prints the counted claim as Theorem~6.9, exactly as the published article does; the displayed formulas inside the proof are numbered sequentially by the wrapper rather than by the article's own section-based scheme. Omitted: the rest of the article, that is its abstract, introduction, the Li--Yau and Aronson--B\'enilan estimates, the preliminary sections, the optimal-control and compactness sections and the appendices, all of which lie outside the source-local core of the counted proof; nothing of the counted statement or of the counted proof is omitted. Class substitution: the source class is the standard \texttt{article} class with the authors' own packages; the wrapper is AMS \texttt{amsart} carrying the authors' own macros, their \texttt{enumitem} lists for the proof steps and their blackboard-bold and indicator macros. Reference handling: no retained passage contains a citation; the article's bibliography exists only in the authors' file \texttt{references.bib}, which is kept verbatim in the eprint directory and is not redistributed, and no reference entry is transcribed or invented. No mathematics is written, repaired or re-derived here.
\setcounter{section}{6}
\setcounter{theorem}{8}
\begin{theorem}\label{thm:new_proof_hls}
Assume that $\rho\ge 0$ is an absolutely continuous measure with finite logarithmic moment such that $\mathcal{F}[\rho]<+\infty$. Take $V\ge0$ such that $K_V=\int_{\R^2}e^{-V}\diff x<+\infty$ and
$\int_{\R^2}\rho V\diff x<+\infty$.
\begin{enumerate}
    \item If the mass $M$ of $\rho$ satisfies $M<8\pi$, then there exists a constant $C=C(M,K_V)$ such that
    $$
        \int_{\R^2}\rho\log\rho\diff x\leq C \left(\mathcal{F}[\rho]+\int V\diff\rho \right).
    $$
    \item If $M=8\pi$, since $\rho$ is not a  Dirac mass,  there exist $\varepsilon>0$, and $R>0$ such that for all $x\in\R^2$ we have $\rho(
B_{R}(x))\leq 8\pi-\varepsilon$; take a constant $\kappa$ such that  $2 - \f{\varepsilon}{4\pi}<\kappa<2$. Then, there exists a constant $C=C(\kappa)$  such that
$$
\left(1-\frac{\kappa}{2}\right)\int_{\R^2}\rho\log\rho\diff x +  C(\kappa)(1+\log R)\le \mathcal{F}[\rho].
$$
\end{enumerate}
\end{theorem}

\begin{proof}
We divide the proof in several steps.
\begin{enumeratesteps}
\step[Setting the inequality]
We first exploit the inequality 
$$
ab\le a\log a + e^{b-1}, 
$$
which holds for all $a,b \in \R$ with $a\ge 0$ and it follows directly from the Legendre's duality. Therefore, if we choose $a=\rho$ and $b=\f{1}{\kappa}u-V(x)$ we obtain
\begin{equation}\label{eq:Inequality for entropy}
    \int_{\R^2}\rho u\diff x \le \kappa\int_{\R^2}\rho\log\rho +\kappa e^{-1}\int_{\R^2}e^{\f{u}{\kappa}-V}\diff x+\kappa\int V\diff\rho. 
\end{equation}
In order to estimate the second term of the right-hand side, we observe that, in both the cases of our statement, for every $x$ there exists a value $R(x)\in [R,+\infty]$ such that $\rho(
B_{R}(x))= 8\pi-\varepsilon$ and that it satisfies $R(x)\ge R$. In the first case, we choose $\ve=8\pi-M>0$ and we have $R(x)=+\infty$ for every $x$, while in the second case we have $R(x)<+\infty$ for every $x$.

Let us take into account the assumptions on $R$ from the statement of the proposition. Let us also take advantage of the Jensen's inequality. Then, we can compute in order to obtain that  
\begin{align*}
    e^{\f{u(x)}{\kappa}}&=\exp\left(\int_{\R^2}\f{1}{2\pi\kappa}\log\left(\f{1}{|x-y|}\right)\rho(y)\diff y\right) \\
    & = \exp\left(\int_{\R^2}\log\left(\f{1}{|x-y|^{\f{1}{2\pi\kappa}}}\right)\rho(y)\diff y\right)\\
    & = \exp\left(\int_{B_{R(x)}(x)}\log\left(\f{1}{|x-y|^{\f{1}{2\pi\kappa}}}\right)\rho(y)\diff y+ \int_{(B_{R(x)}(x))^c}\log\left(\f{1}{|x-y|^{\f{1}{2\pi\kappa}}}\right)\rho(y)\diff y\right)\\
    & = \underbrace{\exp\left(\int_{B_{R(x)}(x)}\log\left(\f{1}{|x-y|^{\f{1}{2\pi\kappa}}}\right)\rho(y)\diff y\right)}_{\eqqcolon A} \, \underbrace{\exp\left( \int_{(B_{R(x)}(x))^c}\log\left(\f{1}{|x-y|^{\f{1}{2\pi\kappa}}}\right)\rho(y)\diff y\right)}_{\eqqcolon B}.
\end{align*}
In the following we bound the exponential factors $A$ and $B$. We observe that in the first case, i.e. $R(x) = + \infty$, we have $B_{R(x)}^c=\emptyset$ and $B=1$.

\step[A bound on $A$ and the case of subcritical mass]
Let us start with the bound on $A$. Let us recall that $\rho(B_{R(x)}(x))=8\pi-\varepsilon$. Therefore, we can use  Jensen's inequality in order to recover that 
\begin{align*}
A&= \exp\left(\int_{B_{R(x)}(x)}\log\left(\f{1}{|x-y|^{\f{1}{2\pi\kappa}}}\right)\rho(y)\diff y\right)\\
&= \exp\left(\int_{B_{R(x)}(x)}\log\left(\f{1}{|x-y|^{\f{\rho \left( B_{R(x)}(x) \right)}{2\pi\kappa}}}\right)\f{\rho(y)}{\rho(B_{R(x)}(x))}\diff y\right)\\
&\le \f{1}{8\pi-\varepsilon}\int_{B_{R(x)}(x)}\f{1}{|x-y|^{\f{8\pi-\varepsilon}{2\pi\kappa}}}\rho(y)\diff y.
\end{align*}
We choose $\kappa$ such that $2 - \f{\varepsilon}{4\pi}<\kappa<2$ so that we have $s \coloneqq \f{8\pi-\varepsilon}{2\pi\kappa} < 2$.

Let us now concentrate on the first case, i.e. subcritical mass $M<8\pi$. We then have to estimate
$$
    \int_{\R^2}e^{\f{u}{\kappa}-V}\diff x\leq \f{1}{M}\int_{\R^2}e^{-V(x)}\diff x\int_{\R^2}\f{1}{|x-y|^s}\diff\rho(y).
$$
If we change the order of the integration in the right-hand side  we obtain
$$
\int_{\R^2}e^{\f{u}{\kappa}-V}\diff x\leq \f{1}{M}\int_{\R^2}\diff\rho(y)\int_{\R^2}e^{-V(x)}\f{1}{|x-y|^s}\diff x.
$$
Let us also recall that we have that
\begin{eqnarray*}
\int_{\R^2}e^{-V(x)}\f{1}{|x-y|^s}\diff x&=&\int_{B_1(y)}e^{-V(x)}\f{1}{|x-y|^s}\diff x+\int_{B_1(y)^c}e^{-V(x)}\f{1}{|x-y|^s}\diff x\\
&\leq&\int_{B_1(y)}\f{1}{|x-y|^s}\diff x
   +\int_{\R^2}e^{-V(x)}\diff x.
\end{eqnarray*}
where the two integrals can both be bounded by a constant only depending on $s$ (hence on $\kappa$) and on $K_V$. 
Therefore, in the first part of the statement ($M<8\pi$ and $R(x) = + \infty$), we obtain
$$
\int_{\R^2}\rho u\diff x \le \kappa\int_{\R^2}\rho\log\rho +\kappa\int V\diff\rho+C, 
$$
and the statement follows by adding $(2-\kappa)\int \rho\log\rho$ to both sides,
using $\mathcal F[\rho]=\int\rho\log\rho-\frac12\int\rho u$.
\step[Critical mass, a bound on $B$ and a new bound on $A$]
We now move on to the second case, i.e. $M=8\pi$, where we also have to take into account the second factor. In this case we choose $V=0$. This requires a different estimate for both $A$ and $B$.

We look at $B$ and we note that the mass on the complementary of $B_{R(x)}(x)$ is $\varepsilon$. Thus, 
\begin{align*}
    B \le  \exp\left( \int_{(B_{R(x)}(x))^c}\log\left(\f{1}{R^{\f{1}{2\pi\kappa}}}\right)\rho(y)\diff y\right) \le \exp\left( \log\left(\f{1}{R(x)^{\f{1}{2\pi\kappa}}}\right) \varepsilon\right)  = \frac{1}{R(x)^{\frac{\varepsilon}{2\pi\kappa}}}.
\end{align*}
%where $C$ is a constant depending on $R$.
%
Observe that we have $\frac{\varepsilon}{2\pi\kappa}=\frac 4\kappa-s$. Furthermore, up to changing the value of $\varepsilon$, we can assume $\kappa>0$. Moreover, since $R(x)\ge R$ we observe that for $y\in B_{R(x)}(x)$, we have $R(x)\ge \max\left(R,|x-y|\right)$. We can now compute the second term on the right-hand side of \eqref{eq:Inequality for entropy} and we obtain 
\begin{align*}
\int_{\R^2}e^{\f{u}{\kappa}}\diff x 
&\lesssim \int_{\R^2} \frac{1}{R(x)^{\frac{4}{\kappa}-s}}\int_{B_{R(x)}(x)}\f{1}{|x-y|^s}\rho(y)\diff y\diff x\\
&=\int_{\R^2}\rho(y)\diff y\int_{\R^2}\frac{1}{R(x)^{\frac{4}{\kappa}-s}}\f{1}{|x-y|^s}\rho(y)\ind_{|x-y|\le R(x)}\diff x \\
&\leq \int_{\R^2}\rho(y)\diff y\int_{\R^2} \frac{1}{\max\left(R^{\frac{4}{\kappa}-s}, |x-y|^{\frac{4}{\kappa}-s}\right)}\f{1}{|x-y|^s}\diff x\\
&=\int_{\R^2}\rho(y)\diff y\int_{B_{R}(y)}\frac{1}{R^{\frac{4}{\kappa}-s}}\f{1}{|x-y|^s}\rho(y)\diff x  + \int_{\R^2}\rho(y)\diff\int_{ B_{R}(y)^c}\f{1}{|x-y|^{\frac{4}{\kappa}}}\diff x \\
&= C(\kappa) R^{2-\frac{4}{\kappa}}.
\end{align*}
In the last inequality, we used the fact that the total mass of $\rho$ is $8\pi$ and we explicitly computed the integrals in $x$.
%
%
\iffalse Also, now we choose $\kappa>2 - \f{\varepsilon}{4\pi}$ so that $s\coloneqq\f{8\pi-\varepsilon}{2\pi\kappa}<2.$ Also up to changing $\varepsilon$ we assume $\kappa>0$. 

Now we can compute the second term on the right-hand side and we obtain 
\begin{align*}
\int_{\R^2}e^{-V(x)}e^{\f{u}{\kappa}}\diff x &\lesssim \int_{\R^2} e^{-V(x)}\int_{B_{R(x)}(x)}\f{1}{|x-y|^s}\rho(y)\diff y\diff x\\
&= \int_{\R^2}\int_{B_{R(x)}(0)} e^{-V(x)}\rho(x+y)\f{1}{|y|^s}\diff y\diff x\\
&= \int_{\R^2}\int_{B_{R(x)}(0)} e^{-V(x-y)}\rho(x)\f{1}{|y|^s}\diff y\diff x\\
&=  \int_{\R^2}\int_{|y|\le R} e^{-V(x-y)}\rho(x)\f{1}{|y|^s}\diff y\diff x + \int_{\R^2}\int_{R\le |y|\le R(x)} e^{-V(x-y)}\rho(x)\f{1}{|y|^s}\diff y\diff x\\
&\eqqcolon I_{1}+ I_{2}.
\end{align*}
Since the mass of $\rho$ is $8\pi$,  $e^{-V(x-y)}\le 1$ and $s<2$ we deduce 
$$
I_{1}\le C. 
$$
We estimate the second term as
\begin{align*}
|I_2| &\le \int_{\R^2}\int_{R\le |y|\le R(x)} e^{-V(x-y)}\rho(x)\f{1}{R^s}\diff y\diff x  \\
&\lesssim \int_{\R^2}\int_{\R^2} e^{-V(x-y)}\rho(x)\diff y\diff x\\
&\lesssim \int_{\R^2}\int_{\R^2} e^{-V(y)}\rho(x)\diff y\diff x\\
&\lesssim 1.
\end{align*}

\fi
%
Hence, the conclusion of this argument is that we have shown that for $2 - \f{\varepsilon}{4\pi}<\kappa<2$ there exists a  constant $C(\kappa)$ such that
$$
\int_{\R^2}\rho u\diff x \le \kappa\int_{\R^2}\rho\log\rho +C(\kappa)R^{2-\frac{4}{\kappa}}. 
$$
We can optimize this result. Let us consider the scaling $\rho_{\lambda}(x)=\lambda^{2}\rho(\lambda x)$ which preserves the mass. Then, a similar argument can be performed for $\rho_{\lambda}$ and we obtain that
$$
    \int_{\R^2}\rho_{\lambda} u_{\lambda}\diff x \le \kappa\int_{\R^2}\rho_\lambda\log\rho_\lambda +C(\kappa)\left(\frac{R}{\lambda}\right)^{2-\frac{4}{\kappa}}. 
$$
We notice that we have
\begin{equation*}
    \int_{\R^2}\rho_{\lambda} u_{\lambda} = \int_{\R^2}\rho u + 32\pi \log\lambda\quad \mbox{ and }\quad\int_{\R^2}\rho_{\lambda}\log\rho_{\lambda}= \int_{\R^2}\rho\log\rho + 16\pi \log\lambda.  
\end{equation*}
Therefore, we have
$$
    \int_{\R^2}\rho u\diff x \le \kappa\int_{\R^2}\rho\log\rho +C(\kappa)\left(\frac{R}{\lambda}\right)^{2-\frac{4}{\kappa}}+(\kappa-2) 16\pi\log\lambda. 
$$
We can optimize  in $\lambda$. We choose $\lambda=R$ and this yields that there exists a new constant $C(\kappa)$ such that 
$$
    \int_{\R^2}\rho u\diff x \le \kappa\int_{\R^2}\rho\log\rho +C(\kappa)(1+\log R)
$$
Once more, we use the definition of $\mathcal F$ and we obtain the desired result. \qedhere
\end{enumeratesteps}
\end{proof}

\end{document}
